%!TEX root = ../../main.tex
\section{지평과 동료 평가기에 대한 안정성}
\label{sec:value-horizon}

결과 시점 규칙의 지평과 평가기를 바꾸었을 때 무엇이 유지되는지는 두 질문으로 나뉜다. 라벨이 유지되는가, 그리고 라벨이 바뀌어도 동결 예측의 상대적 순서가 유지되는가. 두 질문의 답은 조건별로 다르므로 따로 보고한다(표 \ref{tab:horizon-strata-T}, \ref{tab:horizon-strata-S}; 보완 실험, 15.16 노출 이후의 진단).

\paragraph{지평.} 지평 60 / 90 / 120초의 $\SVI$를 비교하면 뒤집힘 비율은 T 전체에서 h60--h90 \hzTAllFlipSixty{}, h90--h120 \hzTAllFlipHundredTwenty{}이고 S에서 \hzSAllFlipSixty{}과 \hzSAllFlipHundredTwenty{}이다. 두 지평이 같은 결과 시점을 공유하는 비율은 h60--h90에서 T \hzTAllSameEp{}, S \hzSAllSameEp{}이고 h90--h120에서 T \hzTSameEpNinety{}이며, 같은 결과 시점을 공유하는 행에서 뒤집힘은 0이다. 뒤집힘은 결과 시점이 다른 행(T \hzTAllFlipDiffEp{}, S \hzSAllFlipDiffEp{})과 작은 변화량($|\dV| < 0.05$; T \hzTSmallFlipSixty{}, S \hzSSmallFlipSixty{})에 집중된다. 지평 간 일치의 상당 부분은 같은 다음 킬에서 잘리는 구간의 동일성에서 오므로, 이 수치는 독립적인 세 지평의 합리성이 아니라 결과 시점 규칙의 안정성을 기술한다.

\paragraph{라벨 이식.} 동결 $q$와 $\PTflex$의 예측을 고정한 채 라벨만 $Y_{h60}$으로 바꾸면 경기 가중 $\dBrier$는 T 전체에서 \hzTAllDbAlt{}(주 라벨 \hzTAllDbMain{}), S 전체에서 \hzSAllDbAlt{}(\hzSAllDbMain{})이다. 범위별 값은 표에 있으며, 작은 변화량 범위에서 $\dBrier$의 절댓값이 가장 작다.

%% GENERATED by thesis/tools/gen_supp.py -- do not edit
\begin{table}[!t]
  \centering
  \caption{지평 층화와 h60 라벨 이식 (한타 T, 15.16 TEST)}
  \label{tab:horizon-strata-T}
  \scriptsize

  \begin{tabular}{@{}lrrrrrrrr@{}}
    \toprule
    범위 & $n$ & 뒤집힘 60$\leftrightarrow$90 & 뒤집힘 90$\leftrightarrow$120 & 같은 종점 60$\leftrightarrow$90 & 뒤집힘$\mid$같은 종점 & 뒤집힘$\mid$다른 종점 & $\Delta$Brier ($Y_{90}$) & $\Delta$Brier ($Y_{60}$) \\
    \midrule
    전체 & 32{,}981 & 0.015 & 0.007 & 0.594 & 0.000 & 0.037 & -0.00373 & -0.00368 \\
    $B_{40}$ & 5{,}423 & 0.014 & 0.006 & 0.553 & 0.000 & 0.031 & -0.00214 & -0.00219 \\
    $|\Delta\hat V|<0.05$ & 13{,}151 & 0.031 & 0.013 & 0.652 & 0.000 & 0.090 & -0.00136 & -0.00107 \\
    같은 프레임 & 2{,}224 & 0.000 & 0.000 & 1.000 & 0.000 & --- & -0.00084 & -0.00084 \\
    새 프레임 & 30{,}757 & 0.016 & 0.007 & 0.565 & 0.000 & 0.037 & -0.00390 & -0.00384 \\
    $t\in[2,10)$ & 2{,}265 & 0.015 & 0.008 & 0.611 & 0.000 & 0.037 & -0.00285 & -0.00219 \\
    $t\in[10,20)$ & 11{,}858 & 0.019 & 0.008 & 0.557 & 0.000 & 0.044 & -0.00525 & -0.00545 \\
    $t\in[20,30)$ & 15{,}446 & 0.013 & 0.006 & 0.588 & 0.000 & 0.031 & -0.00222 & -0.00215 \\
    $t\ge30$ & 3{,}412 & 0.011 & 0.006 & 0.744 & 0.000 & 0.045 & -0.00388 & -0.00360 \\
    \bottomrule
  \end{tabular}
  \tabnote{뒤집힘 = 두 지평의 SVI 라벨이 다른 행의 비율(비가중); 같은 종점 = 두 지평의 결과 시점이 같은 행의 비율. 마지막 두 열은 동결 $q$, $\mathrm{PT}_{\mathrm{flex}}$ 예측을 고정한 채 라벨만 $Y_{h90}$ 또는 $Y_{h60}$으로 바꿔 계산한 경기 가중 $\Delta$Brier(구간 없음). 라벨 이식 진단이며 지평을 다시 설계한 결과가 아니다.}
\end{table}

%% GENERATED by thesis/tools/gen_supp.py -- do not edit
\begin{table}[!t]
  \centering
  \caption{지평 층화와 h60 라벨 이식 (소규모 교전 S, 15.16 TEST)}
  \label{tab:horizon-strata-S}
  \scriptsize

  \begin{tabular}{@{}lrrrrrrrr@{}}
    \toprule
    범위 & $n$ & 뒤집힘 60$\leftrightarrow$90 & 뒤집힘 90$\leftrightarrow$120 & 같은 종점 60$\leftrightarrow$90 & 뒤집힘$\mid$같은 종점 & 뒤집힘$\mid$다른 종점 & $\Delta$Brier ($Y_{90}$) & $\Delta$Brier ($Y_{60}$) \\
    \midrule
    전체 & 101{,}205 & 0.018 & 0.006 & 0.684 & 0.000 & 0.057 & -0.00335 & -0.00326 \\
    $B_{40}$ & 31{,}675 & 0.018 & 0.006 & 0.672 & 0.000 & 0.055 & -0.00210 & -0.00185 \\
    $|\Delta\hat V|<0.05$ & 40{,}468 & 0.037 & 0.013 & 0.698 & 0.000 & 0.122 & -0.00028 & -0.00021 \\
    같은 프레임 & 14{,}026 & 0.000 & 0.000 & 1.000 & 0.000 & --- & -0.00120 & -0.00120 \\
    새 프레임 & 87{,}179 & 0.021 & 0.007 & 0.633 & 0.000 & 0.057 & -0.00351 & -0.00339 \\
    $t\in[2,10)$ & 50{,}444 & 0.020 & 0.007 & 0.678 & 0.000 & 0.061 & -0.00306 & -0.00288 \\
    $t\in[10,20)$ & 33{,}415 & 0.017 & 0.006 & 0.696 & 0.000 & 0.056 & -0.00423 & -0.00422 \\
    $t\in[20,30)$ & 14{,}872 & 0.016 & 0.005 & 0.667 & 0.000 & 0.047 & -0.00106 & -0.00102 \\
    $t\ge30$ & 2{,}474 & 0.013 & 0.006 & 0.722 & 0.000 & 0.045 & +0.00082 & +0.00084 \\
    \bottomrule
  \end{tabular}
  \tabnote{뒤집힘 = 두 지평의 SVI 라벨이 다른 행의 비율(비가중); 같은 종점 = 두 지평의 결과 시점이 같은 행의 비율. 마지막 두 열은 동결 $q$, $\mathrm{PT}_{\mathrm{flex}}$ 예측을 고정한 채 라벨만 $Y_{h90}$ 또는 $Y_{h60}$으로 바꿔 계산한 경기 가중 $\Delta$Brier(구간 없음). 라벨 이식 진단이며 지평을 다시 설계한 결과가 아니다.}
\end{table}


\paragraph{동료 평가기.} 같은 입력의 로지스틱 회귀 평가기(\texttt{\peerName{}}; 기존 선정 이력이 있는 후보이며 TEST에서 새로 고르지 않았다)로 같은 예측·결과 시점을 점수화하면 $\dV$의 부호는 T \peerN{} 행 가운데 \peerSignAgree{}에서 fit85와 일치하고, 라벨은 \peerFlipAll{}의 행에서 뒤집힌다(표 \ref{tab:peer-transfer}). 동결 예측을 고정하고 라벨만 동료 평가기의 $\SVI$로 바꾸면($q$ 입력의 $\ppre$는 교체하지 않음) $\dBrier$는 전체에서 \peerAllDbAlt{}(주 라벨 \peerAllDbMain{}), $\Bforty$에서 \peerBfortyDbAlt{}이고, 작은 변화량 범위에서는 \peerSmallDbAlt{}로 부호가 바뀐다. 즉 라벨의 일치성은 인접한 평가기 사이에서 높지만, 예측 비교의 안정성은 변화량의 크기에 따라 다르다. 이 결과는 동료 평가기로 처음부터 학습한 파이프라인의 성능이 아니라 고정 예측의 라벨 이식 진단이며, 그 의존성은 제 \ref{sec:disc-evaluator} 절에서 다룬다.

%% GENERATED by thesis/tools/gen_supp.py -- do not edit
\begin{table}[!t]
  \centering
  \caption{동료 평가기 라벨 이식 (한타 T, 15.16 TEST)}
  \label{tab:peer-transfer}
  \small

  \begin{tabular}{@{}lrrrr@{}}
    \toprule
    범위 & $n$ & 라벨 뒤집힘 & $\Delta$Brier ($Y_{\mathrm{main}}$) & $\Delta$Brier ($Y_{\mathrm{peer}}$) \\
    \midrule
    전체 & 32{,}981 & 0.086 & -0.00373 & -0.00125 \\
    $B_{40}$ & 5{,}423 & 0.072 & -0.00214 & -0.00067 \\
    $|\Delta\hat V|<0.05$ & 13{,}151 & 0.187 & -0.00136 & +0.00355 \\
    같은 프레임 & 2{,}224 & 0.060 & -0.00084 & -0.00107 \\
    새 프레임 & 30{,}757 & 0.088 & -0.00390 & -0.00118 \\
    $t\in[2,10)$ & 2{,}265 & 0.106 & -0.00285 & -0.00287 \\
    $t\in[10,20)$ & 11{,}858 & 0.093 & -0.00525 & -0.00412 \\
    $t\in[20,30)$ & 15{,}446 & 0.081 & -0.00222 & +0.00129 \\
    $t\ge30$ & 3{,}412 & 0.078 & -0.00388 & +0.00186 \\
    \bottomrule
  \end{tabular}
  \tabnote{동료 평가기 = 같은 입력의 로지스틱 회귀(\texttt{A\_LR\_expanded\_fit85}, 기존 선정 이력 있음; TEST에서 새로 고르지 않음). $Y_{\mathrm{peer}}$ = 동료 평가기로 같은 예측·결과 시점을 점수화한 SVI. 동결 $q$, $\mathrm{PT}_{\mathrm{flex}}$ 예측은 고정, $q$ 입력의 $p_{\mathrm{pre}}$도 교체하지 않음. 경기 가중, 구간 없음.}
\end{table}

